\section{Experimental Setup}
\label{sec:experiments}

\subsection{Dataset}
\label{sec:dataset}

The working dataset is CBIS-DDSM~\cite{lee2017cbis}. BUSI~\cite{aldhabyani2020busi} remains a candidate for later comparison. Mass and calcification tables are merged. After that merge, some patients appear in both official train and test splits; those patients are assigned entirely to train to prevent data leakage. As stated in Section~\ref{sec:concepts}, pathology is reduced to a binary label by treating benign and benign-without-callback as a single negative class. This is a simplifying assumption. Table~\ref{tab:dataset} reports the resulting split sizes.

\begin{table}[!t]
\caption{CBIS-DDSM split used in this study (patient-disjoint after leakage fix).}
\label{tab:dataset}
\centering
\small
\begin{tabular}{@{}lccc@{}}
\toprule
Split & Rows & Patients & Malignant rate \\
\midrule
Train & 2462 & 1060 & 0.414 \\
Validation & 462 & 188 & 0.392 \\
Test & 644 & 318 & 0.399 \\
\bottomrule
\end{tabular}
\end{table}

\subsection{Neural Model Training}
\label{sec:neural}

Concept perception uses an ImageNet-pretrained ResNet18~\cite{he2016resnet} shared trunk with four heads (lesion kind
; 2 classes, shape: 3 classes, margin 4 classes, calcification type: 3 classes). CBIS cropped ROI images are already lesion-centered patches of the full mammogram; each ROI is turned into one fixed $224\times 224$ tensor by random resized crop on train and center crop on validation/test. Grayscale inputs are repeated to three channels and ImageNet-normalized. Geometric augmentation (horizontal flip, small rotation, and mild affine perturbation) is applied to training images only. Because malignant cases are the minority ($\approx 41\%$ of train), augmentation is class-balanced: both classes are expanded to $4329$ training images ($8658$ total).

Before DeepProbLog training, the trunk and heads receive a supervised concept warm-up: eight epochs of class-weighted cross-entropy on the CSV labels for lesion kind, shape, margin, and calcification type (inverse-frequency weights; trunk frozen for the first two warm-up epochs, then fine-tuned at $0.1\times$ the head learning rate). The warmed networks are then fine-tuned for 15 epochs through DeepProbLog ExactEngine at learning rate $3\times 10^{-4}$ and batch size $2$.

A matched end-to-end ResNet18 malignant/benign binary classifier is trained on the same tensors and splits (no DeepProbLog) as Baseline~1. Decision thresholds for both DeepProbLog and the baseline are chosen on validation (including a high-sensitivity target of $0.80$) and applied unchanged to testing.

\subsection{DeepProbLog Configuration}
\label{sec:dplconfig}

Neural predicates supply lesion\_kind, shape, margin, and calc. Symbolic rules gate mass evidence and calcification evidence by lesion kind, as in Section~\ref{sec:knowledge}. The diagnostic query is malignant(X). These rules are a computational working program and are not presented as official BI-RADS criteria. Training jointly updates the concept heads (and, after warmup, the shared trunk) through the probabilistic program.

\subsection{Implemented Architectures}
\label{sec:Implemented Architectures}
Two architectures have been implemented:
\begin{itemize}
\item \textbf{Baseline:} Pure ResNet18 image classifier (same preprocess and splits; no DeepProbLog).

\item \textbf{Proposed architecture:} Shared ResNet18 concepts + DeepProbLog + symbolic medical knowledge.
\end{itemize}

\subsection{Evaluation Metrics}
\label{sec:metrics}

Reported metrics are accuracy, precision, sensitivity, specificity, F1 for both classes and their macro average, and ROC-AUC. The decision threshold on $P(\mathrm{malignant})$ is chosen on validation only. The primary operating point targets validation sensitivity $\geq 0.80$, reflecting the clinical priority of catching malignancies even when that lowers specificity; a more balanced macro-F1 cut-off and the conventional $0.5$ baseline cut-off are also reported. Interpretability---the reason DeepProbLog is used---is assessed qualitatively via reasoning traces on selected true/false positive and negative cases: each trace lists concept probabilities, rule-body strengths, the strongest rule, and the final decision.
